% ----------------------------------------------------------------------------
%                           FUNDAMENTAÇÃO TEÓRICA
% ----------------------------------------------------------------------------

\chapter{FUNDAMENTAÇÃO TEÓRICA}
\label{chap:fundamentacao}

\subsection{Transtorno de Déficit de Atenção com Hiperatividade (TDAH)}

O Transtorno de Déficit de Atenção com Hiperatividade (TDAH) é um transtorno do neurodesenvolvimento de origem neurobiológica, cujos sintomas manifestam-se frequentemente nos primeiros anos de vida e podem persistir até a fase adulta \cite{barkley2015}. De acordo com o Manual Diagnóstico e Estatístico de Transtornos Mentais (DSM-5) da Associação Americana de Psiquiatria \cite{apa2014}, o TDAH caracteriza-se por um padrão persistente de desatenção e/ou hiperatividade-impulsividade que interfere diretamente no funcionamento e no desenvolvimento do indivíduo.

Segundo Mattos \cite{mattos2003}, o transtorno afeta fortemente as funções executivas do cérebro, prejudicando a capacidade de planejamento, a memória de trabalho e a autorregulação. No contexto educacional, alunos com TDAH frequentemente enfrentam desafios para acompanhar o ritmo das atividades propostas. Rohde et al. \cite{rohde2004} enfatizam que essas manifestações não se tratam de incapacidade intelectual ou de falta de vontade, mas de déficits estruturais na regulação atencional e motora.

\subsection{Transtorno Opositivo Desafiador (TOD)}

O Transtorno Opositor Desafiador (TOD) é classificado pelo DSM-5 \cite{apa2014} no grupo dos transtornos disruptivos, do controle de impulsos e da conduta. Ele é descrito como um padrão persistente de humor raivoso e irritável, comportamento questionador, argumentativo e de índole vingativa. Para que o diagnóstico seja estabelecido, tais comportamentos devem ocorrer de forma frequente e persistente, estendendo-se além do que é considerado típico para a idade e o nível de desenvolvimento da criança.

Crianças diagnosticadas com TOD apresentam uma tendência marcante de perder a paciência facilmente, responsabilizar terceiros por seus erros e desafiar ativamente regras ou pedidos feitos por figuras de autoridade, como pais e professores \cite{barkley2015}. É fundamental compreender que o comportamento opositor inerente ao TOD não se resume a uma desobediência comum ou a episódios isolados de frustração, mas reflete uma profunda dificuldade de regulação emocional e uma rigidez cognitiva significativa frente a quebras de expectativa.

\subsection{A Comorbidade entre TDAH e TOD no Contexto Escolar}

Em grande parte dos casos clínicos, o TDAH não se manifesta de maneira isolada. O TOD é uma das comorbidades mais prevalentes em crianças com TDAH. \citeonline{barkley2015} destaca que a coexistência dessas duas condições torna os desafios de manejo clínico e escolar exponencialmente mais complexos. 

Enquanto o TDAH contribui de forma primária com a impulsividade, a agitação motora e a desorganização, o TOD adiciona uma camada de resistência ativa e intencional. No ambiente escolar, essa associação pode resultar em severas desestruturações na dinâmica da sala de aula. O aluno pode não apenas ter dificuldade em iniciar uma tarefa (devido ao TDAH), como também se recusar ativamente a fazê-la e confrontar o professor que a solicitou (devido ao TOD). Consequentemente, esses alunos estão sob alto risco de sofrerem punições frequentes, suspensões e exclusão social, evidenciando a necessidade urgente de abordagens baseadas em evidências em vez de métodos puramente punitivos.

\subsection{O Manejo Comportamental e a Parceria Escola-Família}

Diante das complexidades impostas pelo TDAH e pelo TOD, a literatura médica e pedagógica converge para a necessidade de um tratamento multimodal. Uma das frentes mais importantes dessa abordagem é o manejo comportamental em rede, ou seja, a aplicação de estratégias unificadas e consistentes tanto no ambiente escolar quanto no familiar. Intervenções proativas, como a antecipação e previsibilidade de rotinas, o uso de comandos curtos, a redução de disputas diretas de poder e o reforço positivo imediato, são fundamentais para o desenvolvimento socioemocional desses alunos \cite{mattos2003}.

Para que o manejo comportamental alcance resultados concretos, a consistência é o pilar central. \citeonline{dessen2007} ressaltam que a parceria família-escola atua como um fator de proteção vital para crianças neurodivergentes. Contudo, a ausência de canais de comunicação eficientes rotineiramente provoca o desalinhamento de estratégias --- como a aplicação de uma abordagem rígida na escola contrastando com uma altamente permissiva em casa ---, fato que confunde o aluno e eleva o nível de ansiedade e oposição. Por isso, soluções tecnológicas estruturadas tornam-se ferramentas essenciais para a integração dessa rede de apoio, mitigando ruídos e garantindo que intervenções pedagógicas sejam aplicadas de forma colaborativa e transparente.

\subsection{Contexto Institucional e Normativo: O Projeto TEiA no Instituto Federal de Brasília}

No contexto do Instituto Federal de Brasília (IFB), a efetivação das políticas de educação inclusiva ocorre mediante a atuação estruturada de comissões e projetos institucionais especializados. Destaca-se, nesse cenário, o Projeto TEiA (Formação Continuada em Inclusão e Suporte para Estudantes com TEA e Neuroatípicos), uma iniciativa voltada ao fortalecimento da cultura de inclusão e acessibilidade. A comissão e a equipe do TEiA atuam de forma transversal, envolvendo especialistas em inclusão e operando em parceria direta com a Coordenação de Inclusão (CINC) e as coordenações pedagógicas e de curso do campus \cite{brasil2025ted}.

Os normativos e diretrizes estabelecidos pelo Projeto TEiA definem como prioridade a capacitação dos servidores e o suporte direto às equipes pedagógicas, garantindo o acompanhamento efetivo de estudantes neurodivergentes, o que abrange o manejo comportamental de discentes com TDAH e TOD. Dentre as principais exigências metodológicas e operacionais impostas por essa rede de apoio, sobressai a obrigatoriedade da elaboração, execução e monitoramento dos Planos de Atendimento Educacional Especializado (PAEE) e dos Planos Educacionais Individualizados (PEI). O PEI, em particular, é um instrumento pedagógico central que exige o registro sistemático da rotina diária, documentando as adaptações curriculares e as intervenções táticas aplicadas ao estudante.

A execução dessas diretrizes institucionais, no entanto, esbarra no desafio logístico de unificar os registros e facilitar o compartilhamento dessas informações entre todos os professores e a família. É precisamente nessa lacuna que a proposta de desenvolvimento de um sistema web se justifica, alinhando-se de forma direta aos objetivos do Projeto TEiA.

Ao fornecer uma plataforma centralizada para o registro das ocorrências e para o acompanhamento da rotina diária exigida no PEI, a tecnologia operacionaliza as normativas de inclusão do IFB. Dessa forma, o sistema transcende a função de mero repositório de dados e passa a atuar como um facilitador do fluxo de trabalho da comissão do TEiA, mitigando ruídos de comunicação e garantindo que o monitoramento e a avaliação das práticas inclusivas sejam cumpridos de maneira estruturada, colaborativa e eficiente.

\section{Sistemas Similares}

A fim de mapear o panorama tecnológico atual e identificar as lacunas existentes no mercado educacional e de saúde, torna-se necessário analisar sistemas correlatos já existentes. Essa investigação tem como objetivo compreender como diferentes plataformas abordam o suporte comportamental, a comunicação escolar e a organização de rotinas voltadas ao público neurodivergente, justificando, assim, a proposição da nova solução delineada neste trabalho.

Nesse contexto, são investigados cinco sistemas com diferentes finalidades e abordagens: o \textbf{DBT Coach}, voltado ao acompanhamento clínico e emocional; o \textbf{ClassDojo}, direcionado à gamificação da comunicação entre escola e família; o \textbf{Pictalk Buddy} e o \textbf{Minha Rotina Especial}, focados na organização visual de rotinas e no suporte à previsibilidade das atividades diárias; e a \textbf{Sequência Mágica}, ferramenta educacional voltada ao desenvolvimento da organização lógica e sequencial na educação infantil.

A análise comparativa dessas ferramentas considera aspectos estruturais e funcionais, tais como: objetivo da plataforma, público-alvo, funcionalidades oferecidas, recursos de acessibilidade, estratégias de interação e usabilidade. A partir desse contraste, busca-se compreender de que maneira as soluções já consolidadas podem contribuir para a construção de uma proposta tecnológica mais integrada, acessível e adequada às necessidades táticas de crianças neurodivergentes, de seus responsáveis e dos profissionais da educação.

\vspace{0.3cm}
\subsection{DBT Coach}

O DBT Coach \textit{(Dialectical Behavior Therapy)} é um aplicativo disponível para dispositivos móveis \textit{(Google Play e App Store)} desenvolvido para auxiliar no ensino de habilidades de regulação emocional, conforme informações do seu site oficial \cite{swasth2026}. O software fundamenta-se nos princípios da Terapia Comportamental Dialética e tem como propósito central oferecer suporte clínico para o manejo de crises. Seu público-alvo é composto por pacientes em acompanhamento psicológico e terapeutas que buscam monitorar a evolução do tratamento de forma remota.

O software permite o registro estruturado de estados emocionais e comportamentos por meio da análise funcional (Antecedente, Comportamento e Consequência). Entre suas principais funcionalidades, o aplicativo oferece sugestões de intervenção em tempo real e módulos de treinamento de habilidades. O sistema gera alertas e lembretes baseados no nível de estresse reportado pelo usuário, funcionando como uma ferramenta de suporte imediato para momentos de desregulação.

Em relação à usabilidade, o sistema apresenta uma interface densa e técnica, voltada ao ambiente de saúde mental. Conforme ilustrado na Figura 2.1, a tela principal utiliza um esquema de cores escuras e organiza as funções em um grid de botões para registros rápidos, identificados por ícones para Emoção, Impulso, Comportamento e Sono. A tela de registro clínico exibe uma lista detalhada de comportamentos-alvo, como autolesão e agressão, permitindo que o usuário selecione a intensidade e a frequência de cada ocorrência por meio de seletores táteis.
\begin{figure}[H] 
    \centering
    \caption{Telas principais do aplicativo DBT Coach} % <-- Legenda ÚNICA no topo
    \label{fig:telas_dbt}
    \vspace{0.3cm}
    
    % --- Imagem A ---
    \begin{minipage}[b]{0.46\textwidth}
        \centering
        \includegraphics[height=9cm, keepaspectratio]{dados/figuras/telaDBTCoach_1} 
        \centerline{\small (a) Interface de registros} % <-- Subtítulo A
    \end{minipage}\hfill
    % --- Imagem B ---
    \begin{minipage}[b]{0.46\textwidth}
        \centering
        \includegraphics[height=9cm, keepaspectratio]{dados/figuras/telaDBTCoach_2}
        \centerline{\small (b) Painel clínico} % <-- Subtítulo B
    \end{minipage}
    
   \vspace{0.3cm}
    \begin{flushleft} % <-- Inicia o alinhamento à esquerda
        \footnotesize Fonte: DBT Coach (2026) % <-- Sem o \legend, apenas o texto!
    \end{flushleft}   
\end{figure}

\vspace{0.3cm}
\subsection{ClassDojo}

O ClassDojo é uma plataforma de gestão escolar e comunicação disponível via web e em aplicativos móveis, amplamente utilizada para conectar a escola à família \cite{classdojo2026}. O software tem como propósito central a gamificação do comportamento em sala de aula e a criação de uma comunidade escolar integrada. Seu público-alvo principal são professores da educação básica, alunos e seus respectivos responsáveis.

O software permite a criação de turmas digitais onde cada aluno é representado por um avatar personalizável. As principais funcionalidades incluem o compartilhamento de fotos e vídeos das atividades e um sistema de atribuição de pontos. Os professores podem configurar competências positivas ou pontos a melhorar, emitindo relatórios automáticos sobre o engajamento dos estudantes. O aplicativo ainda permite a troca de mensagens diretas e privadas entre a escola e os pais.

A interface do aplicativo aposta em um design lúdico e colorido para atrair o público infantil. Conforme ilustrado na Figura 2.2, a tela de gestão da turma apresenta avatares em formato de monstrinhos que mudam conforme a pontuação recebida. Ao selecionar um aluno, o sistema abre um menu flutuante com ícones verdes para reforços positivos (como "Trabalho duro") e ícones vermelhos com sinal de menos (-1) para comportamentos que precisam de ajuste (como "Falta de foco"), oferecendo um feedback visual imediato sobre a conduta do estudante.
\begin{figure}[H] 
    \centering
    \caption{Telas principais do aplicativo ClassDojo} % <-- Legenda ÚNICA no topo
    \label{fig:telas_dbt}
    \vspace{0.3cm}
    
    % --- Imagem A ---
    \begin{minipage}[b]{0.46\textwidth}
        \centering
        \includegraphics[height=9cm, keepaspectratio]{dados/figuras/telaClassDojo_1} 
        \centerline{\small (a) Visão geral da turma} % <-- Subtítulo A
    \end{minipage}\hfill
    % --- Imagem B ---
    \begin{minipage}[b]{0.46\textwidth}
        \centering
        \includegraphics[height=9cm, keepaspectratio]{dados/figuras/telaClassDojo_2}
        \centerline{\small (b) Interface de pontuações} % <-- Subtítulo B
    \end{minipage}
    
    \vspace{0.3cm}
    \begin{flushleft} % <-- Inicia o alinhamento à esquerda
    \footnotesize{Fonte: ClassDojo (2026)} % <-- Fonte ÚNICA no final
 \end{flushleft}   
\end{figure}

\vspace{0.2cm}

Contudo, o sistema apresenta falhas metodológicas quando aplicado ao público com TDAH e TOD. O modelo de gamificação baseada em punição pública (perda de pontos visível) pode gerar ansiedade e reforçar o comportamento opositor. Adicionalmente, as avaliações de usuários indicam que o software falha ao não permitir o registro detalhado e privado de ocorrências comportamentais complexas, transformando a comunicação em uma linha do tempo genérica que não orienta os pais sobre como intervir tecnicamente em casa. 

\vspace{0.3cm}
\subsection{Pictalk Buddy}
 
 O Pictalk Buddy é um aplicativo de tecnologia assistiva disponível para smartphones, desenvolvido por uma organização focada em acessibilidade cognitiva \cite{pictalk2026}. O software permite a criação de agendas visuais para auxiliar pessoas com dificuldades na função executiva. Seu público-alvo compreende crianças e adultos com TEA, TDAH e outras condições que se beneficiam da previsibilidade das tarefas diárias.

O software permite a criação passo a passo das rotinas por meio de pictogramas e recursos de voz. As principais funcionalidades envolvem a organização sequencial de atividades, onde o aplicativo gera alertas sonoros que lembram o usuário sobre cada etapa a ser realizada. O sistema também oferece suporte para comunicação alternativa, permitindo que o usuário expresse necessidades básicas através de imagens caso encontre dificuldades na fala verbal.

A usabilidade do Pictalk Buddy é pautada na clareza visual e no minimalismo para evitar a sobrecarga sensorial. Conforme ilustrado na Figura 2.3, a interface é composta por grandes cartões quadrados que contêm ilustrações simples de tarefas cotidianas, como "Escovar os dentes" ou "Lanche". Abaixo de cada imagem, há uma legenda em texto e um botão de áudio; ao ser tocado, o aplicativo narra a instrução em voz alta, permitindo que o usuário avance na rotina de forma autônoma e acompanhe visualmente o progresso do seu dia.
\begin{figure}[H] 
    \centering
    \caption{Telas principais do aplicativo PicTalk} % <-- Legenda ÚNICA no topo
    \label{fig:telas_dbt}
    \vspace{0.3cm}
    
    % --- Imagem A ---
    \begin{minipage}[b]{0.46\textwidth}
        \centering
        \includegraphics[height=9cm, keepaspectratio]{dados/figuras/telaPicTalk_1} 
        \centerline{\small (a) Agenda visual no aplicativo} % <-- Subtítulo A
    \end{minipage}\hfill
    % --- Imagem B ---
    \begin{minipage}[b]{0.46\textwidth}
        \centering
        \includegraphics[height=9cm, keepaspectratio]{dados/figuras/telaPicTalk_2}
        \centerline{\small (b) Organização de agenda de café da manhã} % <-- Subtítulo B
    \end{minipage}
    
  \vspace{0.3cm}
    \begin{flushleft} % <-- Inicia o alinhamento à esquerda
   \footnotesize{Fonte: PicTalk (2026)} % <-- Fonte ÚNICA no final
 \end{flushleft}   
\end{figure}

\vspace{0.2cm}
Apesar de ser eficaz na promoção da autonomia individual, a ferramenta não supre as necessidades de uma gestão escolar compartilhada. Por ser um aplicativo pago e focado no uso pessoal da criança, ele não oferece recursos para que o professor registre comportamentos desafiadores ou alinhe estratégias táticas com os pais. O software atua na prevenção por meio da rotina, mas falha em prover um canal de comunicação e suporte para o manejo comportamental reativo necessário no contexto do TDAH e do TOD.

\vspace{0.3cm}
\subsection{Sequência Mágica}

O sistema web "Sequência Mágica", desenvolvido como pesquisa acadêmica no Instituto Federal de Brasília \cite{santos2025}, é uma plataforma interativa voltada para a organização de rotinas. O propósito principal da ferramenta é auxiliar no desenvolvimento da autonomia e do raciocínio lógico por meio da estruturação de atividades cotidianas. Seu público-alvo direto compreende crianças de 2 a 4 anos diagnosticadas com Transtorno do Espectro Autista (TEA), bem como educadores da primeira infância e pais que auxiliam na mediação dessas atividades.

As principais funcionalidades do sistema baseiam-se na criação interativa de sequências visuais. A plataforma permite que o usuário selecione blocos representativos de ações do dia a dia (como "acordar", "tomar banho" e "escovar os dentes") e os posicione na ordem correta utilizando o recurso de arrastar e soltar (drag and drop). Além disso, o software possui recursos para salvar a rotina montada no armazenamento local do navegador (localStorage), permitindo a visualização, restauração ou exclusão de sequências previamente configuradas e associadas à data atual.

A Figura 2.4 apresenta as telas principais do protótipo do sistema Sequência Mágica. Em relação à usabilidade, o aplicativo adota uma interface voltada para a acessibilidade cognitiva do público infantil, priorizando a clareza visual e a simplicidade de navegação. Como pode ser observado na tela de início (a), há o uso de personagens lúdicos e cores suaves que geram empatia imediata. Já na tela de montagem da sequência (b), a estruturação utiliza uma barra lateral com blocos de rotina diária (como acordar e tomar banho), representados por pictogramas coloridos e legendas simples. Essa disposição limpa e interativa no painel principal permite que a criança acompanhe e estruture o seu próprio dia de forma intuitiva, promovendo a previsibilidade e evitando a sobrecarga sensorial.

\begin{figure}[H]
  \centering
  \caption{Telas principais do protótipo de Sequência Mágica}
  % <-- Legenda ÚNICA no topo
  \label{fig:telas_dbt}
  \vspace{0.3cm}

  % --- Imagem A (Maior largura por ser horizontal) ---
  \begin{minipage}[b]{0.60\textwidth}
    \centering
    \includegraphics[width=\textwidth, keepaspectratio]
    {dados/figuras/telaSequenciaMagica_1}
    \vspace{0.1cm}
    \centerline{\small (a) Tela de início} % <-- Subtítulo A
  \end{minipage}\hfill
  % --- Imagem B (Menor largura por ser vertical) ---
  \begin{minipage}[b]{0.35\textwidth}
    \centering
    \includegraphics[width=\textwidth, keepaspectratio]
    {dados/figuras/telaSequenciaMagica_2}
    \vspace{0.1cm}
    \centerline{\small (b) Tela de montagem da sequência
    dos cards} % <-- Subtítulo B
  \end{minipage}

  \vspace{0.3cm}
  \begin{flushleft} % <-- Inicia o alinhamento à esquerda
  \footnotesize{Fonte: Silva (2025, p. 30, 31).} % <-- Fonte ÚNICA no final
  \end{flushleft}
\end{figure}

\vspace{0.2cm}
Apesar de ser uma iniciativa acadêmica de excelência para a educação infantil, a plataforma apresenta limitações metodológicas e técnicas frente aos objetivos deste projeto. O primeiro obstáculo é estrutural: por utilizar o armazenamento local (localStorage) e não possuir sistema de autenticação ou banco de dados em nuvem, o aplicativo restringe os dados a um único dispositivo, impedindo o compartilhamento de informações integradas entre a escola e a família. Ademais, por focar estritamente na organização preventiva de rotinas para crianças de 2 a 4 anos, a ferramenta não contempla um módulo de comunicação escolar, registro de ocorrências ou suporte tático para o manejo reativo de crises associadas ao TDAH e ao TOD no ambiente educacional.
\vspace{0.3cm}

\section{Análise Comparativa}

Após a análise individual das ferramentas, o Quadro \ref{quadro:comparativo_sistemas} apresenta uma síntese comparativa. Fica evidente que, embora existam soluções robustas para a comunicação alternativa (PicTalk), gamificação (ClassDojo) ou rotinas diárias (Sequência Mágica), há uma lacuna no que diz respeito a um sistema escolar que integre o registro de ocorrências à oferta ativa de estratégias de manejo para professores de alunos com TDAH e TOD, necessidade esta que a solução proposta visa suprir.

\begin{table}[H]
    \centering
    \caption{Comparativo entre os sistemas similares e a solução proposta}
    \label{quadro:comparativo_sistemas}
    \renewcommand{\arraystretch}{1.4} 
    \setlength{\tabcolsep}{3pt}
    \small 

    \begin{tabular}{|>{\raggedright\arraybackslash}p{2.4cm}|>{\raggedright\arraybackslash}p{3.8cm}|>{\raggedright\arraybackslash}p{2.5cm}|>{\centering\arraybackslash}p{2.4cm}|>{\centering\arraybackslash}p{2.0cm}|>{\raggedright\arraybackslash}p{1.8cm}|}
        \hline
        \textbf{Sistema} & \textbf{Foco Principal} & \textbf{Público-Alvo} & \textbf{Comunicação \newline Escola-Família} & \textbf{Sugestões \newline de Manejo} & \textbf{Plataforma} \\ \hline
        
        \textbf{DBT Coach} & Regulação emocional e habilidades de enfrentamento clínico. & Pacientes e terapeutas & Não & Sim (Clínico) & App Mobile \\ \hline
        
        \textbf{ClassDojo} & Gestão de sala de aula, engajamento e gamificação de comportamento. & Professores, pais e alunos & Sim & Não & Web e Mobile \\ \hline
        
        \textbf{PicTalk} & Comunicação Alternativa e Aumentativa (CAA) com pictogramas. & Crianças, pais e terapeutas & Não & Não & App Mobile \\ \hline
        
        \textbf{Sequência Mágica} & Organização de rotinas diárias e acessibilidade cognitiva. & Crianças e pais & Não & Não & App Mobile \\ \hline
        
        \textbf{Sistema Proposto} & Registro comportamental contínuo e suporte estratégico pedagógico. & Educadores (TEiA) e pais \newline (Estudantes com TDAH/TOD) & Sim & Sim & Web Resp. \\ \hline
    \end{tabular}
    
    \vspace{0.2cm}
    \legend{Fonte: Elaborado pela autora (2026)}
\end{table}

\vspace{0.3cm}
\section{Considerações do Capítulo}

Este capítulo estabeleceu os alicerces teóricos e institucionais que fundamentam a pesquisa. Inicialmente, foram delineadas as características clínicas do Transtorno de Déficit de Atenção com Hiperatividade (TDAH) e do Transtorno Opositor Desafiador (TOD), evidenciando a complexidade do manejo comportamental e a necessidade vital de uma parceria integrada e consistente entre a escola e a família. Em seguida, a pesquisa contextualizou essa demanda na realidade normativa do Instituto Federal de Brasília (IFB), pautada pelas diretrizes inclusivas do Projeto TEiA e pela exigência de registros sistemáticos.

Por fim, a análise comparativa de sistemas similares — abrangendo plataformas de comunicação escolar, acessibilidade cognitiva e regulação emocional — comprovou a existência de uma lacuna tecnológica estrutural. A ausência de ferramentas que unam, em um único ambiente, o registro ágil de ocorrências à oferta de estratégias proativas e direcionadas de suporte justifica e sustenta plenamente o desenvolvimento da solução web proposta nos capítulos seguintes.